\subsection{Summary of Chapter 1}

This chapter established the physical and geological foundations required to understand thermoelastic wave propagation in rocks. Geological media were considered as natural solid materials whose behavior is controlled by mineral composition, texture, porosity, pore structure, fractures, bedding, diagenetic history, and in situ stress--temperature conditions. Unlike ideal homogeneous engineering materials, rocks are commonly heterogeneous, anisotropic, and scale-dependent. Their response to thermal and mechanical disturbances therefore reflects both intrinsic material properties and geological structure.

The chapter showed that sandstone, basalt, granite and limestone are distinguished by different physical architectures. Sandstone illustrates the significance of grain supported texture, cementation, pore space and saturation. Basalt represents a dense volcanic rock in which cooling fractures, vesicles, and columnar jointing may strongly influence effective behavior. Granite demonstrates the role of crystalline texture, low primary porosity, jointing, and mineral-scale thermal expansion mismatch. Limestone shows how carbonate composition, dissolution, vugs, fractures, and diagenesis may produce strong heterogeneity. These examples demonstrate that thermoelastic behavior cannot be described by lithology alone, but must be interpreted through the combined effects of mineralogy, structure, porosity, and geological history.

The main physical properties controlling thermoelastic behavior were then introduced. Density determines inertia and therefore influences the acceleration of the medium under stress gradients. Porosity modifies bulk density, stiffness, pore-fluid behavior, heat transfer, and wave propagation. Elastic stiffness, Poisson's ratio, Young's modulus, shear modulus, bulk modulus, and Lamé parameters describe how rocks resist axial deformation, lateral deformation, shear distortion, and volume change. Thermal conductivity, heat capacity, volumetric heat capacity, thermal diffusivity, and the thermal expansion coefficient describe how rocks transmit heat, store thermal energy, smooth temperature disturbances, and deform under temperature change.

Elastic deformation was discussed as the reversible deformation of a solid body under applied stress. In the context of wave propagation, small elastic deformation is especially important because elastic waves are disturbances of displacement, strain, and stress traveling through a deformable medium. The displacement field describes particle motion, strain describes local deformation, and stress describes the internal force state associated with that deformation. Stress gradients produce motion, while the elastic properties of the medium provide the restoring forces that allow waves to propagate. In this framework, P-waves are associated mainly with compressional deformation and volume change, whereas S-waves are associated mainly with shear deformation. Their velocities are controlled by density and elastic moduli, although in real rocks these velocities are also affected by porosity, cracks, saturation, pressure, temperature, and anisotropy.

The chapter also established the thermal part of the thermoelastic problem. Heat conduction was described as a diffusive process governed by thermal conductivity and heat capacity. Elastic waves are a mechanical disturbance with a finite speed . Heat conduction gradually equalizes the temperature differences in the medium . Temperature changes cause thermal expansion or contraction and the deformation is represented by the thermal strain. Where thermal expansion or contraction is restrained by adjacent grains, layers, fractures, joints or the surrounding rock mass, thermal strain results in thermal stress. This mechanism is fundamental to understanding the mechanical effects of heating and cooling of rocks.

The concept of coupled thermoelastic behavior connects these mechanical and thermal processes. Temperature, deformation, stress, and wave motion are not independent in a thermoelastic medium. A temperature perturbation may produce thermal strain, thermal stress, deformation, and elastic motion; at the same time, deformation may be associated with changes in the thermal field, especially in transient processes. This coupling is particularly important in rocks because their internal structure causes temperature gradients, stress concentrations, elastic stiffness, and wave velocities to vary spatially.

Finally, the chapter explained why thermoelastic wave propagation may be direction-dependent. Layering, bedding, foliation, aligned pores, fractures, microcracks, mineral fabrics, fluid saturation, pressure conditions, and temperature gradients can all make physical properties vary with direction. As a result, heat may be conducted more efficiently in some directions than others, and elastic waves may propagate with different velocities, amplitudes, and attenuation depending on their path through the rock. These principles form the physical basis for interpreting directional thermoelastic wave behavior in real geological materials and provide the foundation for the following chapters.